% Auto-generated by src/analysis/latex_tables.py — do not edit by hand.
\begin{table}[htbp]
\centering
\small
\caption{Identification support for the religion vs ethnicity decomposition: cross-county correlation between the 1871 Catholic share and the 1871 Polenpartei vote share, by sub-sample}
\label{tab:cath_polen_identification}
\begin{tabular}{lcccc}
\toprule
Sub-sample & Counties & Mean cath. & Mean Polen. & corr(cath, polen) \\
\midrule
Polish (POS, BRO) & 22 & 65.3 & 60.5 & +0.926 \\
German Catholic (KOL, KOB, TRI, AAC, OPP, MUN) & 63 & 84.9 & 0.0 & n/a \\
Protestant remainder & 254 & 20.7 & 3.2 & +0.301 \\
\midrule
Full panel & 339 & 35.5 & 6.3 & +0.261 \\
\bottomrule
\end{tabular}
\begin{minipage}{\linewidth}
\vspace{0.5em}
\footnotesize \textit{Notes:} Each row reports county-level means and the cross-county correlation between \texttt{cath\_share} (1871 Catholic population share) and \texttt{polen\_share\_1871} (1871 Polenpartei Reichstag vote share) on the relevant sub-sample. The full-panel correlation of $+0.26$ is modest, so the continuous decomposition in Table~\ref{tab:continuous_polish_decomposition} has identifying content. Within the Polish provinces (Posen, Bromberg) the correlation is $+0.93$ -- nearly collinear -- so the decomposition cannot tell apart religion from ethnicity \emph{within} that sub-sample. The German Catholic counties have zero variance in Polenpartei voting (no county in those Regierungsbezirke recorded any Koło Polskie vote), so the within-group correlation is undefined; these counties contribute only to identifying the religion-only coefficient $\beta_1$. The Protestant remainder correlation of $+0.30$ comes from a small number of Danzig and Marienwerder counties with non-zero Polish vote share but low Catholic share. See Figure~\ref{fig:cath_polen_scatter} in the appendix for the corresponding county-level scatter.
\end{minipage}
\end{table}
